\section{The retrocausal worldtube interpretation}
\label{sec:worldtube}

\subsection{Time-symmetry and the arrow of time}
\label{subsec:time-symmetry}

The brane action (\Cref{eq:brane-action}) is time-symmetric.
Its Lorentzian sign structure selects a \emph{kind} of direction --- timelike vs spacelike ---
but the action is invariant under reversal of the timelike coordinate, $t\mapsto -t$.
Neither the kinetic term nor the elastic potential distinguishes past from future.
The substrate therefore has no intrinsic arrow of time: both forward- and backward-in-time
solutions to \Cref{eq:brane-eom} are equally valid, and a configuration is fixed by boundary data
on the past \emph{and} future endcaps of the 4D world-volume --- which is what makes the
retrocausal reading of \Cref{subsec:retrocausal-completion} available.

The empirical arrow we experience --- irreversibility, records of the past but not the future ---
is then not produced by the substrate dynamics. It is the \emph{thermodynamic} arrow: it rests on
a low-entropy boundary condition at one temporal end of the brane, the cosmological ``past
hypothesis'' \citep{Albert2000TimeAndChance,PriceWharton2015Retrocausality} --- the same
low-entropy datum named as the one residual tuning of the ontology
(\Cref{subsec:no-superdeterminism}). On this point the present reading inherits the standard
time-symmetric account, in which the experienced arrow is statistical-mechanical, not dynamical.

\subsection{Entanglement as V-branching worldtubes}
\label{subsec:v-branching}

An entangled pair, on the present reading, is \emph{one} branching 4D world-volume whose worldtube
splits at a V-vertex in the shared past (\Cref{fig:worldtube}).
The branches carry the entangled excitations themselves; for the polarization-entangled photon
pair we commit to (\Cref{subsec:photon-v-worldtube}) these are unconfined carrier waves --- the
photons --- and the only solitons involved are the bound electrons in the detectors, not the
entangled objects.

\begin{figure}[t]
\centering
\begin{tikzpicture}[>={Stealth[length=2.4mm]}, line width=0.8pt]
  % time axis
  \draw[->, gray!70] (-3.7,-1.7) -- (-3.7,2.1) node[left,black]{\footnotesize time};
  % key events
  \coordinate (S)  at ( 0,-1.3);
  \coordinate (DA) at (-2.5, 1.7);
  \coordinate (DB) at ( 2.5, 1.7);
  % the two carrier arms (the branching worldtube)
  \draw[line width=2.4pt, blue!45] (S) -- (DA);
  \draw[line width=2.4pt, blue!45] (S) -- (DB);
  % retrocausal correlation path D_A -> S -> D_B
  \draw[->, red!75, dashed] (DA) .. controls (-1.25,-0.1) .. (S);
  \draw[->, red!75, dashed] (S)  .. controls ( 1.25,-0.1) .. (DB);
  % nodes
  \fill (S) circle (2.3pt);  \node[below] at (S) {\footnotesize source vertex $S$};
  \fill (DA) circle (2.3pt); \node[above] at (DA) {\footnotesize $D_A$ (analyzer $a$)};
  \fill (DB) circle (2.3pt); \node[above] at (DB) {\footnotesize $D_B$ (analyzer $b$)};
  \node[blue!45!black] at (-1.95,0.55) {\footnotesize arm $A$};
  \node[blue!45!black] at ( 1.95,0.55) {\footnotesize arm $B$};
  \node[red!75!black, align=center] at (0,0.55) {\footnotesize correlation path\\[-2pt]\footnotesize (time-symmetric)};
\end{tikzpicture}
\caption{The entangled pair as one branching 4D worldtube. The source vertex $S$ (a nonlinear
matter vertex in the shared past) emits two carrier arms $A$, $B$ that propagate forward to the
detector events $D_A$, $D_B$, where analyzers at settings $a$, $b$ project the polarization. A
setting enters as a \emph{future} boundary condition; because the substrate action is
time-symmetric it propagates backward along that arm to $S$ and forward along the other (dashed),
correlating the outcomes with no spacelike influence. The detectors are spacelike-separated, yet
every step of the dashed path is timelike along a worldline.}
\label{fig:worldtube}
\end{figure}

\paragraph{Mechanism.}
Each branch of the world-volume propagates forward to a measurement event.
The boundary condition imposed by a measurement on one branch is not transmitted
\emph{sideways} through space to the other branch (which would be a spacelike influence,
violating (L)).
Instead, it propagates \emph{backward} along the branch's own worldtube to the V-vertex in the
shared past, and from there \emph{forward} along the other branch to the distant measurement.
The V-vertex is the mediator; no spacelike signal is required.

\paragraph{Locality along each worldline.}
The correlation is:
\begin{itemize}
  \item \emph{local along each worldline}: no influence crosses a spacelike boundary at any
  step of the propagation.
  \item \emph{deterministic given the full 4D configuration}: the correlation is fixed by the
  worldtube geometry and the V-vertex boundary condition.
  \item \emph{Bell-inequality-violating}: the joint measurement context --- both settings $a$ and
  $b$ --- affects the past of both branches via the V-vertex, so the standard factorizability
  assumption of Bell's derivation fails.
\end{itemize}
The 4D world-volume picture promotes ``spooky action at a distance'' to ``ordinary continuity of
one extended elastic object.''

\paragraph{No preferred foliation required.}
Unlike Bohmian mechanics, this picture does not require a preferred spacelike foliation.
The correlation propagates along the worldtube in a covariant, 4D-in-4D sense.
The substrate's emergent effective Lorentz symmetry constrains which foliations are observationally
available to inside observers, but the worldtube itself is a 4D geometric object prior to any
foliation choice.

\subsection{The polarization-entangled photon pair}
\label{subsec:photon-v-worldtube}

The canonical laboratory realization of a V-branching worldtube is the polarization-entangled
photon pair of the entangled-photon Bell tests recognized by the 2022 Nobel Prize --- in
particular the strict-Einstein-locality experiment of Weihs et al.\ \citep{Weihs1998}, built on the
type-II down-conversion source of Kwiat et al.\ \citep{Kwiat1995}, and the loophole-free photon
test of the Vienna group \citep{Giustina2015}.

\paragraph{Apparatus, in substrate terms.}
A pump excitation drives a nonlinear matter vertex --- the down-conversion crystal, a localized
region of bound matter solitons --- in which one carrier wavepacket is converted into two. The two
daughters are unconfined carrier waves, not solitons (\Cref{sec:substrate};
\Cref{subsec:non-probabilistic}). The entangled pair
is therefore a single branching \emph{carrier} world-volume --- one tube that splits at the source
vertex $S$ into two arms, $A$ (toward Alice) and $B$ (toward Bob). The V-vertex of
\Cref{subsec:v-branching} is here the source event $S$; the two measurement endpoints are detector
events $D_A$, $D_B$ at which each arm couples to a bound electron soliton.

\paragraph{The emitted Bell state.}
The vertex coupling is the substrate's geometric three-wave nonlinearity
(\Cref{eq:geom-nonlinear}): a longitudinal pump leg converts into two transverse daughters, the
transverse factor entering through the rotation-invariant norm $|\vecu_\perp|^2$. Its polarization
content is the symmetric contraction $\delta_{ab}$, so the emitted pair is the maximally-entangled
\emph{symmetric} Bell state
\begin{equation}
  \lvert\Phi^+\rangle \;=\; \tfrac{1}{\sqrt 2}\bigl(\lvert H\rangle_A\lvert H\rangle_B
    + \lvert V\rangle_A\lvert V\rangle_B\bigr),
  \label{eq:bellstate}
\end{equation}
a symmetric, type-I-like output. The choice of Bell state carries no loss of generality: a fixed
waveplate at one wing is a local unitary that maps $\lvert\Phi^+\rangle$ to the singlet
$\lvert\Psi^-\rangle$ of a type-II source (as in \citep{Kwiat1995}) and leaves the CHSH value
unchanged, so the substrate's symmetric output and the laboratory state are the same physics up to
a local optical convention. Each arm carries a normalized transverse polarization
$\bm{\epsilon}_A,\bm{\epsilon}_B\in\C^{2}$ (\Cref{eq:carrier});
each detector imposes a \emph{local} boundary condition --- an analyzer at orientation $a$ (Alice)
or $b$ (Bob), and a bound-electron transition rendering the outcome $A_a,B_b\in\{+1,-1\}$, the
discreteness supplied by the detector (\Cref{subsec:non-probabilistic}).

\subsection{How the worldtube reproduces the Born correlations}
\label{subsec:born-correlations}

The branching worldtube reproduces the measured statistics through four pieces, and it matters
which is which. Two of them --- no-signalling and the correlation/Tsirelson bound --- are
\emph{standard quantum-mechanical results}: they follow by the usual Hilbert-space arguments once
the substrate is granted the $\C^{2}$ carrier of \Cref{sec:substrate}, and we use rather than
rederive them. The substrate's own contributions are the other two: it \emph{supplies} that carrier
and emits the Bell state, and it \emph{relocates} the probabilities from a fundamental measure to
classical typicality. Each piece is labeled accordingly below.

\paragraph{No-signalling.}
Bob is ignorant of Alice's outcome, so his local marginal sums over it. Alice's two analyzer-port
states form an orthonormal basis, hence $\sum_{A}\lvert a_A\rangle\langle a_A\rvert=\mathbf{1}$ for
\emph{every} setting $a$, and the sum erases the distant setting:
\begin{equation}
  P(B\mid a,b)=\langle b_B\rvert\,\bigl(\mathrm{Tr}_A\,\lvert\Phi^+\rangle\langle\Phi^+\rvert\bigr)\,\lvert b_B\rangle
  =\langle b_B\rvert\rho_B\rvert b_B\rangle=\tfrac12 ,
  \label{eq:no-signalling}
\end{equation}
independent of $a$ (the one-wing reduced state $\rho_B=\tfrac12\mathbf{1}$ is rotation-invariant).
The setting-dependence carried backward to $S$ therefore lives entirely in the \emph{joint}
correlation, never in a one-wing marginal: Bell nonlocality without signalling. This is the
standard reduced-state no-signalling argument (used here, not a substrate theorem); the substrate's
role is to supply its premises --- the $\C^{2}$ carrier and complete local detection. ``Complete''
means unit detection efficiency, the regime the loophole-free tests reach, not the fair-sampling
assumption those experiments were specifically designed to avoid \citep{Giustina2015,Shalm2015}.

\paragraph{The correlation and Tsirelson's bound.}
A linear polarization at physical angle $\theta$ maps to the Poincar\'e-sphere point at azimuth
$2\theta$ (the geometric phase of polarization \citep{Pancharatnam1956,Berry1984}); the analyzer
observable is $A_a=\bm{s}(a)\!\cdot\!\bm{\sigma}$ with $\bm{s}(a)=(\sin 2a,0,\cos 2a)$. Because the
$\C^{2}$ carrier (\Cref{sec:substrate}) makes the joint history amplitudes genuine Hilbert overlaps, the
correlation is the Poincar\'e inner product
\begin{equation}
  E(a,b)=\langle A_a B_b\rangle = +\cos\!\bigl(2(a-b)\bigr),
  \label{eq:cos-correlation}
\end{equation}
the factor of $2$ being the double angle. At $(a,a',b,b')=(0,\tfrac\pi4,\tfrac\pi8,\tfrac{3\pi}8)$,
\begin{equation}
  \lvert S\rvert=\bigl\lvert E(a,b)-E(a,b')+E(a',b)+E(a',b')\bigr\rvert = 2\sqrt 2 ,
  \label{eq:chsh-tsirelson}
\end{equation}
saturating Tsirelson's bound \citep{Tsirelson1980}. Both the correlation law and the bound are
standard $\C^{2}$ Hilbert-space results, used here rather than rederived; the substrate's
contribution is to \emph{supply} that Hilbert structure --- which is also what \emph{caps} the
correlation at $2\sqrt2$: a setting-coupled box not built from $\C^{2}$ overlaps could reach the
algebraic maximum $S=4$ (a Popescu--Rohrlich box \citep{PopescuRohrlich1994}), whereas a linear
complex carrier admits only genuine overlaps, for which Tsirelson's theorem gives
$\lvert S\rvert\le 2\sqrt2$.

\paragraph{The Born weight.}
The four joint histories carry weights $\lvert\langle a_A b_B\vert\Phi^+\rangle\rvert^2$. We do
\emph{not} derive this Born rule; we offer a \emph{relocation} of it, as the paper's interpretive
proposal, not a theorem. The complex amplitude $\psi=q+\ii p$ is a two-dimensional symplectic phase
space; an equal-a-priori (Liouville) measure on it makes the squared modulus a phase-space area, so
the Born exponent $2$ reads as the symplectic dimension and quantum probability appears as the same
typicality as thermal probability rather than a separate postulate. This is a plausibility argument
with one stated premise --- relaxation of the carrier amplitude to the equilibrium measure, the
ordinary ergodic assumption of statistical mechanics --- and is not a substitute for the structural
Born-rule derivations (Gleason-/Busch-type); establishing it for the substrate dynamics is left
open.

\paragraph{Why the future setting is admissible.}
The analyzer setting $a$ enters only as a future boundary condition at $D_A$; the time-symmetry of
the brane action (\Cref{subsec:time-symmetry}) propagates it backward along $A$ to $S$ and forward
along $B$ to $D_B$, so the joint statistics depend on both $a$ and $b$ while every influence
remains timelike along a worldline. This setting-dependent future datum is what lifts the
correlation above the local bound --- a setting-independent state at $S$ yields only the local
sawtooth (\Cref{subsec:no-superdeterminism}) --- while \eqref{eq:no-signalling} keeps every
one-wing marginal setting-independent, so no signal travels along the backward leg.

\subsection{Open problems}
\label{subsec:bell-open}

\begin{enumerate}
  \item \textbf{V-vertex correlation.}
  A solution of the branching boundary-value problem for the entangled-photon pair
  (\Cref{subsec:photon-v-worldtube}) should exhibit the correlation \eqref{eq:cos-correlation} at
  the two detector endpoints, traceable to continuity through the V-vertex; a reading that
  reproduced it only by adding a spacelike coupling would falsify the picture.

  \item \textbf{Closed-form two-time solution (a strengthening).}
  The consistency result of this paper needs only that the branching structure is available and
  hosts the statistics. Exhibiting an explicit closed-form solution of the two-time branching
  boundary-value problem --- with the chiral-characteristic closure that makes it well-posed, and a
  uniqueness argument --- would strengthen that consistency claim to an explicit construction. It
  is left to future work.

  \item \textbf{Standing assumptions, named.}
  The account rests on two conditions that are standard and not quantum-specific: unit detection
  efficiency (the regime the loophole-free tests reach, not a fair-sampling assumption), under
  which the no-signalling sum \eqref{eq:no-signalling} is exact; and relaxation of the carrier
  amplitude to the equal-a-priori (Liouville) equilibrium measure, the ergodic premise shared with
  classical statistical mechanics, under which the Born relocation holds. Establishing the latter
  for the substrate dynamics is the residual probabilistic debt.
\end{enumerate}